% vim:ts=4:sw=4
%
% Copyright (c) 2008-2009 solvethis
% Copyright (c) 2010-2015 Casper Ti. Vector
% Public domain.
%
% 使用前请先仔细阅读 pkuthss 和 biblatex-caspervector 的文档，
% 特别是其中的 FAQ 部分和用红色强调的部分。
% 两者可在终端/命令提示符中用
%   texdoc pkuthss
%   texdoc biblatex-caspervector
% 调出。

% 采用了自定义的（包括大小写不同于原文件的）字体文件名，
% 并改动 ctex.cfg 等配置文件的用户请自行加入 nofonts 选项；
% 其它用户不用加入 nofonts 选项，加入之后反而会产生错误。
%
% 图书馆要求电子版论文的目录必须为黑色，
% 且某些教务要求打印版论文的文字部分为纯黑色而非灰度打印，
% 【因此最终打印和提交论文前，请将“colorlinks”改为“nocolorlinks”。】
\documentclass[UTF8, nocolorlinks, openany, spacing, pdftoc, pkuspace]{pkuthss}
\linespread{1.5}
% 使用 biblatex 排版参考文献，并规定其格式（详见 biblatex-caspervector 的文档）。
% 这里按照英文文献在前，中文文献在后排序（“sorting = ecnty”）；
% 若需按照中文文献在前，英文文献在后排序，请设置“sorting = centy”；
% 若需按照引用顺序排序，请设置“sorting = none”。
\usepackage[backend = biber, style = caspervector, utf8, sorting = none]{biblatex}
% 提供近似于学校所要求的 Times New Roman / Arial 的字体。
\usepackage[defaultsups]{newtxtext}
\usepackage{newtxmath}
% 产生 originauth.tex 里的 \square。
\usepackage{latexsym}

% 按学校要求设定参考文献列表中的条目之内及之间的距离。
\setlength{\bibitemsep}{3bp}
% 对于 linespread 值的计算过程有兴趣的同学可以参考 pkuthss-extra.sty。
\renewcommand*{\bibfont}{\zihao{5}\linespread{1.27}\selectfont}

% 设定文档的基本信息。
% \ctexset{today=big}
\pkuthssinfo{
	cthesisname = {博士研究生学位论文}, ethesisname = {Doctor Thesis},
	ctitle = {中国科学院近代物理研究所博士论文模版},
	etitle = {Doctor Thesis \LaTeX{} Template for the Institute of Modern Physics of Chineses Academy of Sciences},
	cauthor = {周勇},
	eauthor = {Zhou Yong},
	studentid = {0123456789},
	date = {\today},
	school = {某某学院},
	cmajor = {某某专业}, emajor = {Some Major},
	direction = {某某方向},
	cmentor = {某某教授}, ementor = {Prof.\ Somebody},
	ckeywords = {其一，其二}, ekeywords = {KeywdOne, KeywdTwo}
}
% 载入参考文献数据库（注意不要省略“.bib”）。
\addbibresource{thesis.bib}

% 普通用户可删除此段。
\usepackage{color}
\def\pkuthssffaq{%
	\emph{\textcolor{red}{pkuthss 文档模版最常见问题：}}

	在最终打印和提交论文之前，
	请将 pkuthss 文档类选项中的 %
	\texttt{colorlinks} 替换为 \texttt{nocolorlinks}，
	因为图书馆要求电子版论文的目录必须为黑色，
	且某些教务要求打印版论文的文字部分为纯黑色而非灰度打印。

	\texttt{\string\cite}、\texttt{\string\parencite} %
	和 \texttt{\string\supercite} 三个命令分别产生%
	未格式化的、带方括号的和上标且带方括号的引用标记：%
	\cite{test-en}，\parencite{test-zh}、\supercite{test-en, test-zh}。

	若要避免章末空白页，请在调用 pkuthss 文档类时加入 \texttt{openany} 选项。

	如果编译时不出参考文献，
	请参考 \texttt{texdoc pkuthss}“问题及其解决”一章
	“其它可能存在的问题”一节中关于 biber 的说明。
}

% my packages
\usepackage[absolute, overlay]{textpos} %封面制作

\begin{document}
	%%%%%%%%%%%%%%%%%%%%%%%%%
	%%% 封面,版权,声明,致谢 %%%
	\covermatter
	% 自动生成封面。
	\maketitle
	% 从pdf载入封面
% 	\begin{textblock*}{210mm}(0mm, 0mm)
\begin{figure*}[!hbt]
	\centering
	\includegraphics[width=210mm]{fig/ucas_cover_cn.eps}
\end{figure*}
\end{textblock*}

% 	\null\newpage
% 	\begin{textblock*}{210mm}(0mm, 0mm)
\begin{figure*}[!hbt]
	\centering
	\includegraphics[width=210mm]{fig/cover_en.eps}
\end{figure*}
\end{textblock*}


% 	\null\newpage

	% 版权声明。
	% vim:ts=4:sw=4
%
% Copyright (c) 2008-2009 solvethis
% Copyright (c) 2010-2015 Casper Ti. Vector
% All rights reserved.
%
% Redistribution and use in source and binary forms, with or without
% modification, are permitted provided that the following conditions are
% met:
%
% * Redistributions of source code must retain the above copyright notice,
%   this list of conditions and the following disclaimer.
% * Redistributions in binary form must reproduce the above copyright
%   notice, this list of conditions and the following disclaimer in the
%   documentation and/or other materials provided with the distribution.
% * Neither the name of Peking University nor the names of its contributors
%   may be used to endorse or promote products derived from this software
%   without specific prior written permission.
%
% THIS SOFTWARE IS PROVIDED BY THE COPYRIGHT HOLDERS AND CONTRIBUTORS "AS
% IS" AND ANY EXPRESS OR IMPLIED WARRANTIES, INCLUDING, BUT NOT LIMITED TO,
% THE IMPLIED WARRANTIES OF MERCHANTABILITY AND FITNESS FOR A PARTICULAR
% PURPOSE ARE DISCLAIMED. IN NO EVENT SHALL THE COPYRIGHT HOLDER OR
% CONTRIBUTORS BE LIABLE FOR ANY DIRECT, INDIRECT, INCIDENTAL, SPECIAL,
% EXEMPLARY, OR CONSEQUENTIAL DAMAGES (INCLUDING, BUT NOT LIMITED TO,
% PROCUREMENT OF SUBSTITUTE GOODS OR SERVICES; LOSS OF USE, DATA, OR
% PROFITS; OR BUSINESS INTERRUPTION) HOWEVER CAUSED AND ON ANY THEORY OF
% LIABILITY, WHETHER IN CONTRACT, STRICT LIABILITY, OR TORT (INCLUDING
% NEGLIGENCE OR OTHERWISE) ARISING IN ANY WAY OUT OF THE USE OF THIS
% SOFTWARE, EVEN IF ADVISED OF THE POSSIBILITY OF SUCH DAMAGE.


\chapter*{版权声明}
% \markboth{版权声明}{} 
\thispagestyle{empty} %覆盖\chapter*{}中的调用的\thispagestyle{plain},因为此页不需要页眉页脚

任何收存和保管本论文各种版本的单位和个人，
未经本论文作者同意，不得将本论文转借他人，
亦不得随意复制、抄录、拍照或以任何方式传播。
否则一旦引起有碍作者著作权之问题，将可能承担法律责任。

% 若需排版二维码，请将二维码图片重命名为“barcode”，
% 转为合适的图片格式，并放在当前目录下，然后去掉下面 3 行的注释。
%\vfill\noindent
%\includegraphics[height = 5em]{barcode}


	% 原创性声明和使用授权说明。
	% vim:ts=4:sw=4
%
% Copyright (c) 2008-2009 solvethis
% Copyright (c) 2010-2015 Casper Ti. Vector
% All rights reserved.
%
% Redistribution and use in source and binary forms, with or without
% modification, are permitted provided that the following conditions are
% met:
%
% * Redistributions of source code must retain the above copyright notice,
%   this list of conditions and the following disclaimer.
% * Redistributions in binary form must reproduce the above copyright
%   notice, this list of conditions and the following disclaimer in the
%   documentation and/or other materials provided with the distribution.
% * Neither the name of Peking University nor the names of its contributors
%   may be used to endorse or promote products derived from this software
%   without specific prior written permission.
%
% THIS SOFTWARE IS PROVIDED BY THE COPYRIGHT HOLDERS AND CONTRIBUTORS "AS
% IS" AND ANY EXPRESS OR IMPLIED WARRANTIES, INCLUDING, BUT NOT LIMITED TO,
% THE IMPLIED WARRANTIES OF MERCHANTABILITY AND FITNESS FOR A PARTICULAR
% PURPOSE ARE DISCLAIMED. IN NO EVENT SHALL THE COPYRIGHT HOLDER OR
% CONTRIBUTORS BE LIABLE FOR ANY DIRECT, INDIRECT, INCIDENTAL, SPECIAL,
% EXEMPLARY, OR CONSEQUENTIAL DAMAGES (INCLUDING, BUT NOT LIMITED TO,
% PROCUREMENT OF SUBSTITUTE GOODS OR SERVICES; LOSS OF USE, DATA, OR
% PROFITS; OR BUSINESS INTERRUPTION) HOWEVER CAUSED AND ON ANY THEORY OF
% LIABILITY, WHETHER IN CONTRACT, STRICT LIABILITY, OR TORT (INCLUDING
% NEGLIGENCE OR OTHERWISE) ARISING IN ANY WAY OUT OF THE USE OF THIS
% SOFTWARE, EVEN IF ADVISED OF THE POSSIBILITY OF SUCH DAMAGE.

{
	\CTEXsetup[
		format+ = {\centering}, beforeskip = {40bp}, afterskip = {15bp}
	]{section}

	\chapter*{北京大学学位论文原创性声明和使用授权说明}
	\thispagestyle{empty} %覆盖\chapter*{}中的调用的\thispagestyle{plain},因为此页不需要页眉页脚
	\mbox{}\vspace*{-3em}
	\section*{原创性声明}

	本人郑重声明：
	所呈交的学位论文，是本人在导师的指导下，独立进行研究工作所取得的成果。
	除文中已经注明引用的内容外，
	本论文不含任何其他个人或集体已经发表或撰写过的作品或成果。
	对本文的研究做出重要贡献的个人和集体，均已在文中以明确方式标明。
	本声明的法律结果由本人承担。
	\vskip 1em
	\rightline{%
		论文作者签名：\hspace{5em}%
		日期：\hspace{2em}年\hspace{2em}月\hspace{2em}日%
	}

	\section*{%
		学位论文使用授权说明\\[-0.33em]
		\textmd{\zihao{5}（必须装订在提交学校图书馆的印刷本）}%
	}

	本人完全了解北京大学关于收集、保存、使用学位论文的规定，即：
	\begin{itemize}
		\item 按照学校要求提交学位论文的印刷本和电子版本；
		\item 学校有权保存学位论文的印刷本和电子版，
			并提供目录检索与阅览服务，在校园网上提供服务；
		\item 学校可以采用影印、缩印、数字化或其它复制手段保存论文；
		\item 因某种特殊原因需要延迟发布学位论文电子版，
			授权学校在 $\square$\nobreakspace{}一年 / %
			$\square$\nobreakspace{}两年 / %
			$\square$\nobreakspace{}三年以后在校园网上全文发布。
	\end{itemize}
	\centerline{（保密论文在解密后遵守此规定）}
	\vskip 1em
	\rightline{%
		论文作者签名：\hspace{5em}导师签名：\hspace{5em}%
		日期：\hspace{2em}年\hspace{2em}月\hspace{2em}日%
	}

	% 若需排版二维码，请将二维码图片重命名为“barcode”，
	% 转为合适的图片格式，并放在当前目录下，然后去掉下面 2 行的注释。
	%\vfill\noindent
	%\includegraphics[height = 5em]{barcode}
}


	% 致谢。
	% vim:ts=4:sw=4
% Copyright (c) 2014 Casper Ti. Vector
% Public domain.

\chapter*{\hypertarget{ack}{致谢}}
% \thispagestyle{empty} %覆盖\chapter*{}中的调用的\thispagestyle{plain},因为此页不需要页眉页脚
% \pdfbookmark[1]{致谢}{ack}
% 中文测试文字。
\pkuthssffaq


	
	%%%%%%%%%%%%%%%%%%%%%
	%%% 中英文摘要,目录 %%%
	% 简单页眉页脚.此后内容都加入目录,并双面打印
	\frontmatter
	% 中英文摘要。
	% vim:ts=4:sw=4
% Copyright (c) 2014 Casper Ti. Vector
% Public domain.

\begin{cabstract}
	% 中文测试文字。
	\pkuthssffaq
	电子排版系统的出现给印刷出版业带来了一场革命，利用电子计算机及各种辅助设备，可以完成从文稿、图表的录入、编辑、修改、组版，直至得到各种不同用途、不同质量的输出结果。利用电子排版系统，可以减轻劳动强度，缩短出版周期。

　　目前世界上有许多电子排版系统。这些系统各有特点，也各有自己的适用范围。TeX 就是一种优秀的电子排版系统。

　　TeX 提供了一套功能强大并且十分灵活的排版语言，它多达 900 多条指令，并且 TeX 有宏功能，用户可以不断地定义自己适用的新命令来扩展 TeX 系统的功能。许多人利用 TeX 提供的宏定义功能对 TeX 进行了二次开发，其中比较著名的有美国数学学会推荐的非常适合于数学家使用的 AMS-TeX 以及适合于一般文章、报告、书籍的 LaTeX 系统。

　　TeX 系统是公认的数学公式排得最好的系统。美国数学学会（AMS） 鼓励数学家们使用 TeX 系统向它的期刊投稿。世界上许多一流的出版社如 Kluwer、Addison-Wesley、牛津大学出版社等也利用 TeX 系统出版书籍和期刊。

　　大部分的 TeX 系统都是免费的。Knuth 教授还公开了他的全部源程序。TeX 系统目前已经在数百种计算机系统上得到实现。TeX 系统的排版结果 DVI（DeVice Independent）文件与输出设备无关。DVI 文件可以显示、打印、照排，几乎可以在所有的输出设备上输出。TeX 排版源文件及结果在各种计算机系统上互相兼容。

\end{cabstract}

% \cleardoublepage
\begin{eabstract}
	Test of the English abstract.version 1.3c of this license or (at your option) any later
    version. This version of this license is in
       http://www.latex-project.org/lppl/lppl-1-3c.txt
    and the latest version of this license is in
       
       http://www.latex-project.org/lppl.txt
    and version 1.3 or later is part of all distributions of
    LaTeX version 2005/12/01 or later.
    If LaTeX decides to introduce at a point in the document where a  is specified, then no white space is produced. To ensure that white space is produced even at points in the document where page breaking takes place, one should replartain features of LaTeX relating to blank spaces and paragraph indentation which will improve the appearance of the final document. Experienced users of LaTeX will improve the appearance of their documents if they bear these remarks in mind.

First note that, as a general rule, you should never put a blank space after a left parenthesis or before a right parenthesis. If you were to put a blank space in these places, then you run the risk that LaTeX might start a new line immediately after the left parenthesis or before the right parenthesis, leaving the parenthesis marooned at the beginning or end of a line.

LaTeX has its own rules for deciding the lengths of blank spaces. For instance, LaTeX will put an extra amount of space after a full stop if it considers that the full stop marks the end of a sentence.

The rule adopted by LaTeX is to regard a period (full stop) as the end of a sentence if it is preceded by a lowercase letter. If the period is preceded by an uppercase letter then LaTeX assumes that it is not a full stop but follows the initials of somebody's name.

This works very well in most cases. However LaTeX occasionally gets things wrong. This happens with a number of common abbreviations (as in `Mn `etc.'), and, in particular, in the names of journals given in abbreviated form (e.g., `Proc.

oc.'). The way to overcome this problem is to put a backslash before the blank space in question. Thus we should type


LaTeX determines itself how to break up a paragraph into lines, and will occasionally hyphenate long words where this is desirable. However it is sometimes necessary to tell LaTeX not to break at a particular blank space. The special character used for this purpose is ~. It represents a blank space at which LaTeX is not allowed to break between lines. It is often desirable to use ~ in names where the forenames are represented by initials. Thus to obtain `W. R. Hamilton' it is best to type W.~R.~Hamilton. It is also desirable in phrases like `Example 7' and `the length l of the rod', obtained by typing


\end{eabstract}


	% 自动生成目录。	
	\tableofcontents

	%%%%%%%%%%%%%%%%%
	%%% 正文各章节 %%%
	\mainmatter
	% 各章节。
	% vim:ts=4:sw=4
% Copyright (c) 2014 Casper Ti. Vector
% Public domain.

\chapter{sina新闻}
一个多月前，周树春正在北京某训练基地对第43届世界技能大赛做准备，突然接到一个北京号码的陌生来电。电话中，一名男子自称是人社部的工作人员，告诉他：今年9月在北京有一个活动邀请他参加。几天后，单位通知他，人社部邀请他9月去北京参加胜利日大阅兵观礼。此时，他才想起之前接的那个神秘电话说的可能是同一件事。
\section{中文测试}
% 中文测试文字。
\pkuthssffaq
\section{新闻内容}
\subsection{第一个新闻内容}
\subsubsection{第二个新闻内容}
一个多月前，周树春正在北京某训练基地对第43届世界技能大赛做准备，突然接到一个北京号码的陌生来电。电话中，一名男子自称是人社部的工作人员，告诉他：今年9月在北京有一个活动邀请他参加。几天后，单位通知他，人社部邀请他9月去北京参加胜利日大阅兵观礼。此时，他才想起之前接的那个神秘电话说的可能是同一件事。

周树春说，他小时候在电视里看到过阅兵式，虽然隔着小小的电视机屏幕，但那整齐划一的队伍，气势浩荡的口号，一排排军人和各式先进武器行进在天安门前，那个震撼的场面给年幼的他留下了深刻的印象。“幸福来得太突然了！”周树春说，他做梦都没想到自己会成为那个到现场观礼的幸运儿。为了能在观礼台上看清楚，他打算买一个高清相机和望远镜。
周树春说，他小时候在电视里看到过阅兵式，虽然隔着小小的电视机屏幕，但那整齐划一的队伍，气势浩荡的口号，一排排军人和各式先进武器行进在天安门前，那个震撼的场面给年幼的他留下了深刻的印象。“幸福来得太突然了！”周树春说，他做梦都没想到自己会成为那个到现场观礼的幸运儿。为了能在观礼台上看清楚，他打算买一个高清相机和望远镜。
% \linespread{0.1}周树春说，他小时候在电视里看到过阅兵式，虽然隔着小小的电视机屏幕，但那整齐划一的队伍，气势浩荡的口号，一排排军人和各式先进武器行进在天安门前，那个震撼的场面给年幼的他留下了深刻的印象。“幸福来得太突然了！”周树春说，他做梦都没想到自己会成为那个到现场观礼的幸运儿。为了能在观礼台上看清楚，他打算买一个高清相机和望远镜。
% \linespread{1.3}周树春说，他小时候在电视里看到过阅兵式，虽然隔着小小的电视机屏幕，但那整齐划一的队伍，气势浩荡的口号，一排排军人和各式先进武器行进在天安门前，那个震撼的场面给年幼的他留下了深刻的印象。“幸福来得太突然了！”周树春说，他做梦都没想到自己会成为那个到现场观礼的幸运儿。为了能在观礼台上看清楚，他打算买一个高清相机和望远镜。
% \linespread{1.5}周树春说，他小时候在电视里看到过阅兵式，虽然隔着小小的电视机屏幕，但那整齐划一的队伍，气势浩荡的口号，一排排军人和各式先进武器行进在天安门前，那个震撼的场面给年幼的他留下了深刻的印象。“幸福来得太突然了！”周树春说，他做梦都没想到自己会成为那个到现场观礼的幸运儿。为了能在观礼台上看清楚，他打算买一个高清相机和望远镜。
周树春说，他小时候在电视里看到过阅兵式，虽然隔着小小的电视机屏幕，但那整齐划一的队伍，气势浩荡的口号，一排排军人和各式先进武器行进在天安门前，那个震撼的场面给年幼的他留下了深刻的印象。“幸福来得太突然了！”周树春说，他做梦都没想到自己会成为那个到现场观礼的幸运儿。为了能在观礼台上看清楚，他打算买一个高清相机和望远镜。
%  \linespread{10.0}周树春说，他小时候在电视里看到过阅兵式，虽然隔着小小的电视机屏幕，但那整齐划一的队伍，气势浩荡的口号，一排排军人和各式先进武器行进在天安门前，那个震撼的场面给年幼的他留下了深刻的印象。“幸福来得太突然了！”周树春说，他做梦都没想到自己会成为那个到现场观礼的幸运儿。为了能在观礼台上看清楚，他打算买一个高清相机和望远镜。
周树春说，他小时候在电视里看到过阅兵式，虽然隔着小小的电视机屏幕，但那整齐划一的队伍，气势浩荡的口号，一排排军人和各式先进武器行进在天安门前，那个震撼的场面给年幼的他留下了深刻的印象。“幸福来得太突然了！”周树春说，他做梦都没想到自己会成为那个到现场观礼的幸运儿。为了能在观礼台上看清楚，他打算买一个高清相机和望远镜。
周树春说，他小时候在电视里看到过阅兵式，虽然隔着小小的电视机屏幕，但那整齐划一的队伍，气势浩荡的口号，一排排军人和各式先进武器行进在天安门前，那个震撼的场面给年幼的他留下了深刻的印象。“幸福来得太突然了！”周树春说，他做梦都没想到自己会成为那个到现场观礼的幸运儿。为了能在观礼台上看清楚，他打算买一个高清相机和望远镜。

周树春说，他小时候在电视里看到过阅兵式，虽然隔着小小的电视机屏幕，但那整齐划一的队伍，气势浩荡的口号，一排排军人和各式先进武器行进在天安门前，那个震撼的场面给年幼的他留下了深刻的印象。“幸福来得太突然了！”周树春说，他做梦都没想到自己会成为那个到现场观礼的幸运儿。为了能在观礼台上看清楚，他打算买一个高清相机和望远镜。
周树春说，他小时候在电视里看到过阅兵式，虽然隔着小小的电视机屏幕，但那整齐划一的队伍，气势浩荡的口号，一排排军人和各式先进武器行进在天安门前，那个震撼的场面给年幼的他留下了深刻的印象。“幸福来得太突然了！”周树春说，他做梦都没想到自己会成为那个到现场观礼的幸运儿。为了能在观礼台上看清楚，他打算买一个高清相机和望远镜。

周树春说，他小时候在电视里看到过阅兵式，虽然隔着小小的电视机屏幕，但那整齐划一的队伍，气势浩荡的口号，一排排军人和各式先进武器行进在天安门前，那个震撼的场面给年幼的他留下了深刻的印象。“幸福来得太突然了！”周树春说，他做梦都没想到自己会成为那个到现场观礼的幸运儿。为了能在观礼台上看清楚，他打算买一个高清相机和望远镜。
周树春说，他小时候在电视里看到过阅兵式，虽然隔着小小的电视机屏幕，但那整齐划一的队伍，气势浩荡的口号，一排排军人和各式先进武器行进在天安门前，那个震撼的场面给年幼的他留下了深刻的印象。“幸福来得太突然了！”周树春说，他做梦都没想到自己会成为那个到现场观礼的幸运儿。为了能在观礼台上看清楚，他打算买一个高清相机和望远镜。
周树春说，他小时候在电视里看到过阅兵式，虽然隔着小小的电视机屏幕，但那整齐划一的队伍，气势浩荡的口号，一排排军人和各式先进武器行进在天安门前，那个震撼的场面给年幼的他留下了深刻的印象。“幸福来得太突然了！”周树春说，他做梦都没想到自己会成为那个到现场观礼的幸运儿。为了能在观礼台上看清楚，他打算买一个高清相机和望远镜。
周树春说，他小时候在电视里看到过阅兵式，虽然隔着小小的电视机屏幕，但那整齐划一的队伍，气势浩荡的口号，一排排军人和各式先进武器行进在天安门前，那个震撼的场面给年幼的他留下了深刻的印象。“幸福来得太突然了！”周树春说，他做梦都没想到自己会成为那个到现场观礼的幸运儿。为了能在观礼台上看清楚，他打算买一个高清相机和望远镜。
周树春说，他小时候在电视里看到过阅兵式，虽然隔着小小的电视机屏幕，但那整齐划一的队伍，气势浩荡的口号，一排排军人和各式先进武器行进在天安门前，那个震撼的场面给年幼的他留下了深刻的印象。“幸福来得太突然了！”周树春说，他做梦都没想到自己会成为那个到现场观礼的幸运儿。为了能在观礼台上看清楚，他打算买一个高清相机和望远镜。
周树春说，他小时候在电视里看到过阅兵式，虽然隔着小小的电视机屏幕，但那整齐划一的队伍，气势浩荡的口号，一排排军人和各式先进武器行进在天安门前，那个震撼的场面给年幼的他留下了深刻的印象。“幸福来得太突然了！”周树春说，他做梦都没想到自己会成为那个到现场观礼的幸运儿。为了能在观礼台上看清楚，他打算买一个高清相机和望远镜。
周树春说，他小时候在电视里看到过阅兵式，虽然隔着小小的电视机屏幕，但那整齐划一的队伍，气势浩荡的口号，一排排军人和各式先进武器行进在天安门前，那个震撼的场面给年幼的他留下了深刻的印象。“幸福来得太突然了！”周树春说，他做梦都没想到自己会成为那个到现场观礼的幸运儿。为了能在观礼台上看清楚，他打算买一个高清相机和望远镜。

\newpage
周树春说，他小时候在电视里看到过阅兵式，虽然隔着小小的电视机屏幕，但那整齐划一的队伍，气势浩荡的口号，一排排军人和各式先进武器行进在天安门前，那个震撼的场面给年幼的他留下了深刻的印象。“幸福来得太突然了！”周树春说，他做梦都没想到自己会成为那个到现场观礼的幸运儿。为了能在观礼台上看清楚，他打算买一个高清相机和望远镜。
周树春说，他小时候在电视里看到过阅兵式，虽然隔着小小的电视机屏幕，但那整齐划一的队伍，气势浩荡的口号，一排排军人和各式先进武器行进在天安门前，那个震撼的场面给年幼的他留下了深刻的印象。“幸福来得太突然了！”周树春说，他做梦都没想到自己会成为那个到现场观礼的幸运儿。为了能在观礼台上看清楚，他打算买一个高清相机和望远镜。
周树春说，他小时候在电视里看到过阅兵式，虽然隔着小小的电视机屏幕，但那整齐划一的队伍，气势浩荡的口号，一排排军人和各式先进武器行进在天安门前，那个震撼的场面给年幼的他留下了深刻的印象。“幸福来得太突然了！”周树春说，他做梦都没想到自己会成为那个到现场观礼的幸运儿。为了能在观礼台上看清楚，他打算买一个高清相机和望远镜。


	% 结论。
	% vim:ts=4:sw=4
% Copyright (c) 2014 Casper Ti. Vector
% Public domain.

\chapter{结论}
% 中文测试文字。
\section{中文测试}
\pkuthssffaq

\section{通知一}

各学院、各位研究生：

为维护研究生考试的公平、公正，严肃考试纪律，保证本学期期末考试工作的顺利进行，现将有关事项通知如下。

一、考试时间及考场安排

研究生期末考试定于2015年1月5日至1月16日进行。公共课由研究生院及开课学院安排，专业课由各学院负责组织安排与实施。公共课考试的具体安排见附件一、附件二及附件三。

二、考试纪律

研究生参加考试必须严格遵守考场纪律，具体要求见附件四。对违反考试纪律、不服从监考教师管理的违纪、作弊考生，将根据情节轻重参照《兰州大学学生违反考试纪律处理办法》的有关规定予以处理。

三、有关要求

各学院要高度重视考试工作，严格组织考务及有关工作，组织我校正式教职工担任监考，并对监考人员严格要求，加强培训。对研究生加强考试纪律要求，严肃考风考纪。

各学院要成立由主管院长负责的考试工作小组，安排对专业课考试的巡考工作。研究生院将组织教学督导专家及工作人员开展巡考工作。

各学院务必通知研究生携带证件按时参加考试。研究生须认真阅读考场规则及考试安排，按要求时间地点准时参加考试，以免耽误考试。

\subsection{再次通知}
各学院、各位研究生：

为维护研究生考试的公平、公正，严肃考试纪律，保证本学期期末考试工作的顺利进行，现将有关事项通知如下。

一、考试时间及考场安排

研究生期末考试定于2015年1月5日至1月16日进行。公共课由研究生院及开课学院安排，专业课由各学院负责组织安排与实施。公共课考试的具体安排见附件一、附件二及附件三。

二、考试纪律

研究生参加考试必须严格遵守考场纪律，具体要求见附件四。对违反考试纪律、不服从监考教师管理的违纪、作弊考生，将根据情节轻重参照《兰州大学学生违反考试纪律处理办法》的有关规定予以处理。

三、有关要求

各学院要高度重视考试工作，严格组织考务及有关工作，组织我校正式教职工担任监考，并对监考人员严格要求，加强培训。对研究生加强考试纪律要求，严肃考风考纪。

各学院要成立由主管院长负责的考试工作小组，安排对专业课考试的巡考工作。研究生院将组织教学督导专家及工作人员开展巡考工作。

\section{通知二}
各学院务必通知研究生携带证件按时参加考试。研究生须认真阅读考场规则及考试安排，按要求时间地点准时参加考试，以免耽误考试。

为切实提高我校研究生尤其是博士研究生的科技论文英文写作水平，学校将按学科开设研究生科技论文英文写作公共选修课。本学期，研究生院特邀请澳大利亚阿德莱德大学Dr. Margaret Cargill授课。其中，博士研究生培训班在部分专业进行试点，具体事宜另行通知。现就青年教师培训班相关事宜通知如下：

1.培训的主要内容为科技论文的结构、文章的科学逻辑、英文写作及审稿要求、文稿的投递、评审与修改等。

2.培训范围为我校各学科青年教师20-30名，培训结束经考核合格后发放结业证书。参加培训者要承担本学科博士研究生科技论文英文写作课程的教学任务。

3.培训为英语授课，请各学院推荐一名青年教师参加培训，外语学院推荐三名。要求被推荐者具有较高的英语水平，能熟练运用英语进行听、说、读、写，有国外留学经历者优先。

4.请各学院推荐教师填写“科技论文英文写作培训班登记表”（附件1），并于5月4日前交至研究生院培养处（贵勤楼531室）。

5. 授课时间为5月18日-22日，授课地点另行通知。同时请参加培训的老师自行准备教材“Writing Scientific Research Articles-strategy and steps （ Margaret Cargill Patrick O’Connor 著，Wiley-Blackwell出版社2009)”（可自网上下载） 。

请各学院认真遴选推荐青年教师参加培训，并按时报送登记表，未尽事宜请与研究生院培养处联系。

联系人：吕伟东

联系电话：8912660


	% special chapter没有章节编号，在目录中显示
	\specialchap{特殊章节}
% 中文测试文字。
\pkuthssffaq
	
	% 正文中的附录部分。
	\appendix
	% 排版参考文献列表。
	\printbibliography[
		% 使“参考文献”出现在目录中；如果同时要使参考文献列表参与章节编号，
		% 可将“bibintoc”改为“bibnumbered”。
		heading = bibintoc,
		% 单独设定排序方案。此设定会局部覆盖之前的全局设置。
		% 注：只有同时使用 2.x 或之后版本的 biblatex 和相应兼容版本的 biber，
		% 才能对每个 \printbibliography 命令采用不同的排序方案，
		% 否则只能在载入 biblatex 宏包时就（全局）指定排序方案。
		% 在这样的情况下，请去掉所有的 sorting 选项，否则可能出错。
		% 此外，biblatex 3.0 中 \printbibliography 的 sorting 选项失效，
		% 详见 biblatex-caspervector 的文档。
		sorting = ecnty
	]
	% 各附录。
	% vim:ts=4:sw=4
% Copyright (c) 2014 Casper Ti. Vector
% Public domain.

\chapter{附件}
% 中文测试文字。
\pkuthssffaq



	% 以下为正文之后的部分，默认不进行章节编号。
	\backmatter
	
\end{document}

